\subsection{Reading Notes:Transitions to Two Hadron States From QCD}
This is Felipe's PhD thesis. It discusses a mathematical formalism to ascertain the structure of hadrons using LQCD, and uses this formalism to perform a numerical calculation for a certain type of hadron transition. \\\\
- There have been many recently discovered hadrons, most of which are resonances, that need to be understood theoretically. \\\\
- QCD is asymptotically free at high energy: this means that at high energy the QCD coupling constant $\alpha_s$ approaches zero, which allows for perturbative calculations as higher order terms in $\alpha_s$ are negligible. 
- QCD is non-perturbative at low energy: at low energy $\alpha_s ~ 1$, so we cannot perform any perturbative calculations. This regime is what causes quarks to be confined within hadrons. \\\\
- The fact that QCD is non-perturbative at low energy renders most of our tools for understanding and predicting the behaviors of particles inapplicable, and we still do not have a proof for why quark confinement arises, or how thee hadron spectrum arises from quarks and gluons.\\\\
- We classify mesons and baryons as bosons and fermions, respectively. Note: I believe they are further specified as quark-anti quark and three quark formations, but I'm not certain. Hadrons that fall outside of these constructions are called exotic. \\\\
- Understanding these states could lead to a better understanding of non-exotic hadrons, as well as understanding why the spectrum of hadrons arises. \\
- The substructure of hadrons was first experimentally discovered in Stanford in the 1950s, where electrons were scattered off protons. Since then, we've developed other experimental techniques to study the substructure, including electron positron colliders.\\\\ 
- most hadrons are only detectable in a  collider by resonances in energy distributions of their decay products, as their lifetimes are so short that they cannot reach the detector directly.\\\\
- What's the ``path integral of QCD"?  LQCD is a numerical approach to doing this path integral. \\\\
- 